\documentclass[11pt,a4paper]{article}

% ---- Packages ----
\usepackage[utf8]{inputenc}
\usepackage[T1]{fontenc}
\usepackage{amsmath,amssymb,amsthm}
\usepackage{mathtools}
\usepackage[margin=1in,headheight=14pt]{geometry}
\usepackage{hyperref}
\usepackage{enumitem}
\usepackage{xcolor}
\usepackage{tcolorbox}
\usepackage{fancyhdr}
\usepackage{booktabs}

% ---- Hyperref ----
\hypersetup{
    colorlinks=true,
    linkcolor=red,
    citecolor=blue,
    urlcolor=blue
}

% ---- Header ----
\pagestyle{fancy}
\fancyhf{}
\fancyhead[L]{\textit{Lesson 7: Hilbert Spaces and Operators}}
\fancyhead[R]{\thepage}
\renewcommand{\headrulewidth}{0.4pt}

% ---- Theorem Environments ----
\theoremstyle{definition}
\newtheorem{definition}{Definition}[section]
\newtheorem{example}{Example}[section]
\newtheorem{exercise}{Exercise}[section]

\theoremstyle{plain}
\newtheorem{theorem}{Theorem}[section]
\newtheorem{lemma}[theorem]{Lemma}
\newtheorem{proposition}[theorem]{Proposition}
\newtheorem{corollary}[theorem]{Corollary}

\theoremstyle{remark}
\newtheorem{remark}{Remark}[section]

% ---- Colored Boxes ----
\tcbuselibrary{theorems,skins,breakable}

\newtcolorbox{keyresult}[1][]{
    colback=green!5!white,
    colframe=green!50!black,
    fonttitle=\bfseries,
    title=#1,
    breakable
}

\newtcolorbox{rlconnection}[1][]{
    colback=blue!5!white,
    colframe=blue!50!black,
    fonttitle=\bfseries,
    title=#1,
    breakable
}

\newtcolorbox{warning}[1][]{
    colback=red!5!white,
    colframe=red!50!black,
    fonttitle=\bfseries,
    title=#1,
    breakable
}

% ---- Operators ----
\newcommand{\R}{\mathbb{R}}
\newcommand{\N}{\mathbb{N}}
\newcommand{\Q}{\mathbb{Q}}
\newcommand{\Z}{\mathbb{Z}}
\newcommand{\C}{\mathbb{C}}
\newcommand{\Prob}{\mathbb{P}}
\newcommand{\E}{\mathbb{E}}
\newcommand{\indicator}{\mathbf{1}}
\DeclareMathOperator{\interior}{int}
\DeclareMathOperator{\cl}{cl}
\DeclareMathOperator{\spn}{span}
\DeclareMathOperator{\conv}{conv}
\DeclareMathOperator{\dom}{dom}
\DeclareMathOperator{\range}{range}
\DeclareMathOperator{\nullspace}{null}
\DeclareMathOperator{\diam}{diam}

% ---- Title ----
\title{Lesson 19: Inner Product Spaces and Hilbert Spaces}
\author{Curriculum: A Rigorous Learning Path to Multi-Agent Deep RL}
\date{}

\begin{document}
\maketitle
\tableofcontents
\newpage

\setcounter{section}{0}

\section{Introduction}\label{sec:introduction}

The passage from normed spaces to inner product spaces provides the crucial
ingredient of \emph{geometry}: an inner product defines angles, orthogonality,
and projections. These geometric tools are not merely elegant abstractions---they
are the backbone of least-squares methods, kernel machines, and the projection
operators used in temporal-difference learning (LSTD, LSPE).

Complete inner product spaces---\emph{Hilbert spaces}---enjoy the richest
structure among Banach spaces. The Projection Theorem, the Riesz Representation
Theorem, and the spectral theory of operators on Hilbert spaces all depend on
this completeness.

Beyond Hilbert spaces, the four fundamental theorems of functional analysis
(Hahn--Banach, Uniform Boundedness, Open Mapping, Closed Graph) are tools
of immense power for proving existence and continuity results in abstract spaces.

Finally, fixed point theorems translate the abstract structure of function spaces
into concrete algorithmic guarantees:

\begin{itemize}[leftmargin=2em]
    \item The \textbf{Banach Contraction Mapping Theorem} guarantees that
          value iteration converges: the Bellman operator is a
          $\gamma$-contraction on the space of bounded value functions.
    \item The \textbf{Brouwer and Schauder Fixed Point Theorems} guarantee
          existence of fixed points for continuous maps on compact convex sets.
    \item The \textbf{Kakutani Fixed Point Theorem} extends Brouwer to
          set-valued maps, providing the standard proof that every finite
          game has a Nash equilibrium.
\end{itemize}

\begin{rlconnection}[Why This Lesson Matters for RL/ML]
\begin{enumerate}[leftmargin=2em]
    \item \textbf{Projections and least-squares}: Temporal-difference methods
          with linear function approximation (LSTD) project the Bellman
          equation onto a subspace. The Projection Theorem makes this rigorous.
    \item \textbf{Kernel methods}: Reproducing kernel Hilbert spaces (RKHS)
          are Hilbert spaces of functions. The Riesz Representation Theorem
          is the starting point for representer theorems.
    \item \textbf{Value iteration convergence}: The Banach Contraction Mapping
          Theorem, applied to the Bellman operator, gives both existence of
          the optimal value function and the geometric convergence rate
          $\|V_n - V^*\| \leq \gamma^n \|V_0 - V^*\|$.
    \item \textbf{Nash equilibrium existence}: The Kakutani Fixed Point Theorem,
          applied to the best-response correspondence, proves existence of
          Nash equilibria in finite games---the theoretical anchor of MARL.
    \item \textbf{Uniform Boundedness}: Ensures that pointwise-bounded
          families of operators are uniformly bounded, a key tool in proving
          convergence of approximation schemes.
\end{enumerate}
\end{rlconnection}

\medskip
\noindent\textbf{Prerequisites.} Familiarity with metric spaces, normed spaces,
Banach spaces, and bounded linear operators as covered in Lesson~6. Measure
theory from Lessons~1--2 is used in $L^2$ examples.

\medskip
\noindent\textbf{Outline.}
Section~\ref{sec:inner-product} develops inner product spaces and proves the
Cauchy--Schwarz inequality.
Section~\ref{sec:hilbert} introduces Hilbert spaces, orthogonality, and
orthogonal decomposition.
Section~\ref{sec:projection} proves the Projection Theorem and studies
projection operators.
Section~\ref{sec:orthonormal} covers orthonormal systems, Bessel's inequality,
and Parseval's identity.
Section~\ref{sec:riesz} proves the Riesz Representation Theorem.
Section~\ref{sec:fundamental} presents the four fundamental theorems of
functional analysis with complete proofs.
Section~\ref{sec:fixed-point} proves the Banach, Brouwer, Schauder, and
Kakutani fixed point theorems.
Section~\ref{sec:exercises} provides exercises.
Section~\ref{sec:summary} summarizes the key results.


%% ============================================================================
%% SECTION 2: INNER PRODUCT SPACES
%% ============================================================================

\section{Inner Product Spaces}\label{sec:inner-product}

\begin{definition}[Inner Product Space]\label{def:inner-product}
Let $X$ be a vector space over $\mathbb{F} \in \{\R, \C\}$. An
\emph{inner product} on $X$ is a map
$\langle \cdot, \cdot \rangle \colon X \times X \to \mathbb{F}$ satisfying,
for all $x, y, z \in X$ and $\alpha \in \mathbb{F}$:
\begin{enumerate}[label=(\roman*)]
    \item \textbf{Linearity in the first argument:}
          $\langle \alpha x + y, z \rangle = \alpha \langle x, z \rangle
          + \langle y, z \rangle$.
    \item \textbf{Conjugate symmetry:}
          $\langle x, y \rangle = \overline{\langle y, x \rangle}$.
    \item \textbf{Positive definiteness:}
          $\langle x, x \rangle \geq 0$, with equality if and only if $x = 0$.
\end{enumerate}
A vector space equipped with an inner product is called an
\emph{inner product space} (or \emph{pre-Hilbert space}).
\end{definition}

\begin{remark}\label{rem:convention}
We adopt the convention that the inner product is linear in the first argument
and conjugate-linear in the second. Some texts use the opposite convention.
In the real case, conjugate symmetry reduces to symmetry:
$\langle x, y \rangle = \langle y, x \rangle$.
\end{remark}

\begin{example}[Standard Inner Product Spaces]\label{ex:standard-ip}
\begin{enumerate}[label=(\alph*)]
    \item \textbf{Euclidean space $\R^n$:}
          $\langle x, y \rangle = \sum_{i=1}^n x_i y_i$.
    \item \textbf{Sequence space $\ell^2$:}
          $\langle x, y \rangle = \sum_{i=1}^\infty x_i \overline{y_i}$
          for $x = (x_i), y = (y_i) \in \ell^2$.
    \item \textbf{Function space $L^2(\mu)$:}
          $\langle f, g \rangle = \int f \,\overline{g}\, d\mu$
          for $f, g \in L^2(\mu)$.
    \item \textbf{Matrix space $\R^{m \times n}$:}
          $\langle A, B \rangle = \mathrm{tr}(B^\top A)$, the Frobenius
          inner product.
\end{enumerate}
\end{example}

\subsection{The Cauchy--Schwarz Inequality}

The single most important inequality in all of analysis.

\begin{theorem}[Cauchy--Schwarz Inequality]\label{thm:cauchy-schwarz}
Let $X$ be an inner product space. For all $x, y \in X$,
\[
    |\langle x, y \rangle| \leq \|x\| \cdot \|y\|,
\]
where $\|x\| = \langle x, x \rangle^{1/2}$. Equality holds if and only if
$x$ and $y$ are linearly dependent.
\end{theorem}

\begin{proof}
If $y = 0$, both sides are zero and the inequality holds trivially.
Assume $y \neq 0$. For any $\alpha \in \mathbb{F}$,
\[
    0 \leq \langle x - \alpha y, x - \alpha y \rangle
    = \|x\|^2 - \alpha \langle y, x \rangle
      - \overline{\alpha} \langle x, y \rangle
      + |\alpha|^2 \|y\|^2.
\]
Choose $\alpha = \langle x, y \rangle / \|y\|^2$. Then
$\overline{\alpha} = \overline{\langle x, y \rangle} / \|y\|^2
= \langle y, x \rangle / \|y\|^2$, and
\begin{align*}
    0 &\leq \|x\|^2
        - \frac{\langle x, y \rangle}{\|y\|^2} \langle y, x \rangle
        - \frac{\langle y, x \rangle}{\|y\|^2} \langle x, y \rangle
        + \frac{|\langle x, y \rangle|^2}{\|y\|^4} \|y\|^2 \\
      &= \|x\|^2
        - \frac{|\langle x, y \rangle|^2}{\|y\|^2}
        - \frac{|\langle x, y \rangle|^2}{\|y\|^2}
        + \frac{|\langle x, y \rangle|^2}{\|y\|^2} \\
      &= \|x\|^2 - \frac{|\langle x, y \rangle|^2}{\|y\|^2}.
\end{align*}
Rearranging gives $|\langle x, y \rangle|^2 \leq \|x\|^2 \|y\|^2$, and
taking square roots yields the result.

For the equality condition: if $x = \beta y$ for some $\beta \in \mathbb{F}$,
then $|\langle x, y \rangle| = |\beta| \|y\|^2 = \|x\| \|y\|$.
Conversely, equality in the above derivation means
$\langle x - \alpha y, x - \alpha y \rangle = 0$, hence $x = \alpha y$.
\end{proof}

\begin{corollary}[Triangle Inequality]\label{cor:triangle}
$\|x + y\| \leq \|x\| + \|y\|$ for all $x, y \in X$.
\end{corollary}

\begin{proof}
$\|x + y\|^2 = \|x\|^2 + 2\,\mathrm{Re}\,\langle x, y \rangle + \|y\|^2
\leq \|x\|^2 + 2\|x\|\|y\| + \|y\|^2 = (\|x\| + \|y\|)^2$.
\end{proof}

Thus $\|x\| = \langle x, x \rangle^{1/2}$ defines a norm on $X$: the
\emph{induced norm}. Every inner product space is a normed space.

\begin{theorem}[Parallelogram Law]\label{thm:parallelogram}
In an inner product space,
\[
    \|x + y\|^2 + \|x - y\|^2 = 2(\|x\|^2 + \|y\|^2)
    \quad \text{for all } x, y \in X.
\]
\end{theorem}

\begin{proof}
Expand both sides:
\begin{align*}
    \|x + y\|^2 &= \|x\|^2 + 2\,\mathrm{Re}\,\langle x, y \rangle + \|y\|^2, \\
    \|x - y\|^2 &= \|x\|^2 - 2\,\mathrm{Re}\,\langle x, y \rangle + \|y\|^2.
\end{align*}
Adding these gives the result.
\end{proof}

\begin{remark}\label{rem:parallelogram-converse}
The converse holds: a norm satisfying the parallelogram law arises from an inner
product (recovered via the polarization identity). This means, for example,
that $\ell^p$ for $p \neq 2$ is \emph{not} an inner product space.
\end{remark}

\begin{proposition}[Polarization Identity]\label{prop:polarization}
In a real inner product space,
\[
    \langle x, y \rangle
    = \frac{1}{4}\bigl(\|x+y\|^2 - \|x-y\|^2\bigr).
\]
In a complex inner product space,
\[
    \langle x, y \rangle
    = \frac{1}{4}\sum_{k=0}^{3} i^k \|x + i^k y\|^2.
\]
\end{proposition}

\begin{proof}
For the real case, expand $\|x+y\|^2 - \|x-y\|^2 = 4\,\langle x, y \rangle$
using the same computation as in the parallelogram law proof.

For the complex case, let $\omega = i$ and compute:
\begin{align*}
    \sum_{k=0}^{3} i^k \|x + i^k y\|^2
    &= \sum_{k=0}^{3} i^k \bigl(\|x\|^2 + i^k \langle y, x \rangle
       + \overline{i^k} \langle x, y \rangle + \|y\|^2\bigr).
\end{align*}
Since $\sum_{k=0}^{3} i^k = 0$ and $\sum_{k=0}^{3} i^{2k} = 0$,
the $\|x\|^2$ and $\|y\|^2$ terms vanish. We have
$\sum_{k=0}^{3} i^k \cdot \overline{i^k} = \sum_{k=0}^{3} |i^k|^2 = 4$
and $\sum_{k=0}^{3} i^k \cdot i^k = \sum_{k=0}^{3} i^{2k} = 0$.
Therefore the sum equals $4\,\langle x, y \rangle$.
\end{proof}

\begin{proposition}[Continuity of the Inner Product]\label{prop:ip-continuity}
The inner product is jointly continuous: if $x_n \to x$ and $y_n \to y$ in an
inner product space, then $\langle x_n, y_n \rangle \to \langle x, y \rangle$.
\end{proposition}

\begin{proof}
$|\langle x_n, y_n \rangle - \langle x, y \rangle|
\leq |\langle x_n - x, y_n \rangle| + |\langle x, y_n - y \rangle|
\leq \|x_n - x\| \|y_n\| + \|x\| \|y_n - y\|$.
Since convergent sequences are bounded, $\|y_n\| \leq M$ for some $M$,
so the right side tends to zero.
\end{proof}


%% ============================================================================
%% SECTION 3: HILBERT SPACES
%% ============================================================================

\section{Hilbert Spaces}\label{sec:hilbert}

\begin{definition}[Hilbert Space]\label{def:hilbert-space}
A \emph{Hilbert space} is an inner product space that is complete with respect
to the norm induced by its inner product. That is, every Cauchy sequence
converges in the norm $\|x\| = \langle x, x \rangle^{1/2}$.
\end{definition}

\begin{example}[Standard Hilbert Spaces]\label{ex:hilbert-spaces}
\begin{enumerate}[label=(\alph*)]
    \item $\R^n$ with the Euclidean inner product is a Hilbert space (finite-dimensional
          normed spaces are always complete).
    \item $\ell^2 = \{(x_i)_{i=1}^\infty : \sum |x_i|^2 < \infty\}$ is a Hilbert space.
    \item $L^2(\mu) = \{f : \int |f|^2 \, d\mu < \infty\} / {\sim}$ (equivalence
          classes modulo a.e.\ equality) is a Hilbert space by the Riesz--Fischer theorem.
\end{enumerate}
\end{example}

\begin{warning}[Non-Example]
The space $C[0,1]$ with the $L^2$ inner product
$\langle f, g \rangle = \int_0^1 f(t)g(t)\, dt$ is an inner product space
but \emph{not} a Hilbert space: one can construct Cauchy sequences of
continuous functions whose limit in $L^2$ is not continuous.
\end{warning}

\subsection{Orthogonality}

\begin{definition}[Orthogonality]\label{def:orthogonality}
Two elements $x, y$ in an inner product space are \emph{orthogonal}, written
$x \perp y$, if $\langle x, y \rangle = 0$. A set $S \subset X$ is
\emph{orthogonal} if any two distinct elements of $S$ are orthogonal.
\end{definition}

\begin{theorem}[Pythagorean Theorem]\label{thm:pythagorean}
If $x_1, \ldots, x_n$ are pairwise orthogonal, then
\[
    \Bigl\|\sum_{k=1}^n x_k\Bigr\|^2 = \sum_{k=1}^n \|x_k\|^2.
\]
\end{theorem}

\begin{proof}
$\bigl\|\sum_k x_k\bigr\|^2 = \sum_k \sum_j \langle x_k, x_j \rangle
= \sum_k \langle x_k, x_k \rangle = \sum_k \|x_k\|^2$,
since $\langle x_k, x_j \rangle = 0$ for $k \neq j$.
\end{proof}

\begin{definition}[Orthogonal Complement]\label{def:orthogonal-complement}
For a subset $S \subset X$, the \emph{orthogonal complement} of $S$ is
\[
    S^\perp = \{x \in X : \langle x, s \rangle = 0 \text{ for all } s \in S\}.
\]
\end{definition}

\begin{proposition}[Properties of Orthogonal Complements]\label{prop:orth-complement}
Let $S, T$ be subsets of an inner product space $X$. Then:
\begin{enumerate}[label=(\roman*)]
    \item $S^\perp$ is a closed subspace of $X$.
    \item $S \subset T \implies T^\perp \subset S^\perp$.
    \item $S \subset (S^\perp)^\perp$.
    \item $S^\perp = (\overline{\spn(S)})^\perp$.
\end{enumerate}
\end{proposition}

\begin{proof}
\begin{enumerate}[label=(\roman*)]
    \item Linearity of the inner product in the first argument shows $S^\perp$
          is a subspace. Continuity of the inner product
          (Proposition~\ref{prop:ip-continuity}) shows it is closed:
          if $x_n \in S^\perp$ and $x_n \to x$, then for every $s \in S$,
          $\langle x, s \rangle = \lim_n \langle x_n, s \rangle = 0$.
    \item If $x \in T^\perp$, then $\langle x, t \rangle = 0$ for all $t \in T
          \supset S$, so $x \in S^\perp$.
    \item If $s \in S$ and $x \in S^\perp$, then $\langle s, x \rangle = 0$,
          so $s \in (S^\perp)^\perp$.
    \item One inclusion: $S \subset \overline{\spn(S)}$ implies
          $(\overline{\spn(S)})^\perp \subset S^\perp$.
          Other inclusion: if $x \in S^\perp$, then $\langle x, s \rangle = 0$
          for all $s \in S$. By linearity, $\langle x, y \rangle = 0$ for all
          $y \in \spn(S)$. By continuity, $\langle x, y \rangle = 0$ for all
          $y \in \overline{\spn(S)}$. \qedhere
\end{enumerate}
\end{proof}


%% ============================================================================
%% SECTION 4: PROJECTION THEOREM
%% ============================================================================


\end{document}